\section{Admissibility-Constrained Hierarchical RAL-MoE-ROM}
\label{sec:hierarchical_ral_moe}

The frozen specialists are assembled by an outer sparse mechanism that
separates three decisions: regime preference, topological legality, and
dynamical feasibility.  Router probabilities alone do not authorize a
specialist to participate in a prediction.  The resulting hierarchy uses
Top-1 routing in regime cores, permits an output-level Top-2 mechanism only
on a certified adjacent overlap, and fails closed when no candidate satisfies
its native contract.



\subsection{Trajectory-level routing and adjacency topology}
\label{subsec:trajectory_router_topology}

The outer E2 router receives only the standardized Reynolds number,
\begin{equation}
  \widetilde{\mu}
  =
  \frac{\mu-\mu_{\mathrm{R}}}{\sigma_{\mathrm{R}}},
  \qquad
  \boldsymbol{\pi}(\mu)
  =
  \operatorname{softmax}
  \bigl(\mathcal{R}_{E2}(\widetilde{\mu})\bigr)
  =
  [\pi_S,\pi_H,\pi_P]^{\mathsf T}.
  \label{eq:e2_outer_router}
\end{equation}
One routing decision is made for a complete fixed-$\mu$ trajectory.  The
outer router is therefore a trajectory-level, temporal-consistent regime
router, rather than a per-step switching policy.  Because the available
training data do not constitute certified continuous startup or
parameter-ramp transitions, the router is not interpreted as a learned
physical $S\!\to\!H\!\to\!P$ transition model.

The only legal regime adjacency is
\begin{equation}
  A_{\mathrm{adj}}
  =
  \begin{bmatrix}
    1&1&0\\
    1&1&1\\
    0&1&1
  \end{bmatrix},
  \label{eq:adjacency_matrix}
\end{equation}
using the class order $(S,H,P)$.  For an anchor regime $k$, the topological
candidate set is
\begin{equation}
  \mathcal{C}_{\mathrm{top}}(k)
  =
  \{j\in\mathcal{R}:A_{\mathrm{adj},kj}=1\}.
  \label{eq:topological_candidate_set}
\end{equation}
This removes a direct $S$--$P$ pair and excludes Top-3 activation before any
field prediction is evaluated.

\subsection{Dynamical admissibility certification}
\label{subsec:admissibility_certification}

For specialist $j$, initial physical history $\mathcal{H}$, Reynolds number
$\mu$, and registered rollout horizon $K$, define the binary admissibility
indicator
\begin{equation}
  m_j(\mathcal{H},\mu,K)\in\{0,1\}.
  \label{eq:admissibility_indicator}
\end{equation}
The value $m_j=1$ requires all of the following:
\begin{enumerate}
  \item every predicted modal state and reconstructed field is finite;
  \item the registered horizon contains no numerically divergent window;
  \item the local POD, scaler, history, integrator, Pressure--Poisson,
        and pressure-gauge contracts are satisfied;
  \item independent validation trajectories support stable native rollout
        over the same development domain and horizon; and
  \item inference uses neither test truth, a future endpoint, a future state,
        nor an uncertified cross-chart input.
\end{enumerate}
Given router weights $\pi_j$, topology and admissibility produce
\begin{equation}
  \widetilde{\pi}_j
  =
  \frac{
    \mathbf{1}_{\{j\in\mathcal{C}_{\mathrm{top}}\}}\,
    m_j\,\pi_j
  }{
    \sum_{k\in\mathcal{R}}
    \mathbf{1}_{\{k\in\mathcal{C}_{\mathrm{top}}\}}\,
    m_k\,\pi_k+\varepsilon
  }.
  \label{eq:masked_router_weights}
\end{equation}
If two adjacent candidates are admissible, a legal Top-2 comparison is
formed and the subsequent T2-C gate is applied.  If only one is admissible,
the hierarchy degenerates deterministically to that Top-1 specialist.  If
none is admissible, the system returns a fail-closed status instead of
forcing a numerical prediction.  Admissibility is thus a prerequisite for
capability comparison, not a loss term that can be compensated by router
confidence.



\subsection{Boundary-specific hierarchical assembly}
\label{subsec:boundary_specific_assembly}

\paragraph{Regime-core prediction}
In a certified core domain, only its native specialist is executed:
\begin{equation}
  \widehat{\boldsymbol{x}}
  =
  \begin{cases}
    \widehat{\boldsymbol{x}}_{S},& \text{steady core},\\
    \widehat{\boldsymbol{x}}_{H},& \text{Hopf core},\\
    \widehat{\boldsymbol{x}}_{P},& \text{periodic core}.
  \end{cases}
  \label{eq:core_top1_prediction}
\end{equation}
Each prediction is generated in its local chart and reconstructed with its
own decoder.

\paragraph{Top-2 convex correction gate (T2-C)}
For any adjacent regime pair, a lightweight gate is trained to blend the
outputs of the two frozen specialists.  The gate is evaluated once per query
window of $K$ steps and produces a convex weight that remains constant
throughout the window.

The candidate ordering is fixed per boundary: for the $S$--$H$ pair,
Candidate~1 is the Steady specialist and Candidate~2 is the Hopf specialist;
for the $H$--$P$ pair, Candidate~1 is the Periodic specialist and
Candidate~2 is the Hopf specialist.  The E2 probabilities are first
renormalized on the legal pair,
\begin{equation}
  \overline{\pi}_1
  =
  \frac{\pi_1}{\pi_1+\pi_2},
  \qquad
  \overline{\pi}_2=1-\overline{\pi}_1,
  \qquad
  \ell^{\mathrm{E2}}
  =
  \log\frac{\overline{\pi}_1+\varepsilon}
  {\overline{\pi}_2+\varepsilon},
  \label{eq:t2c_pair_logit}
\end{equation}
where $(\pi_1,\pi_2)=(\pi_S,\pi_H)$ for the $S$--$H$ boundary and
$(\pi_1,\pi_2)=(\pi_P,\pi_H)$ for the $H$--$P$ boundary.




\paragraph{Chart-independent physical-history descriptors.}

To provide the overlap-fusion model with information that is independent of
the regime-local POD coordinates, we construct a compact descriptor vector
directly from the three most recent physical states.  Let
\[
  \left\{
    \bigl(
      \boldsymbol{u}_{n-j},
      p_{n-j}^{\circ},
      t_{n-j}
    \bigr)
  \right\}_{j=0}^{2}
\]
denote the current state and the two preceding states, all represented on the
common finite-volume mesh and using the pressure gauge defined in
\eqref{eq:pressure_gauge}.  The two most recent time increments are
\begin{equation}
  \Delta t_n
  =
  t_n-t_{n-1},
  \qquad
  \Delta t_{n-1}
  =
  t_{n-1}-t_{n-2}.
  \label{eq:descriptor_time_increments}
\end{equation}

Let $U_{\mathrm{ref}}$, $P_{\mathrm{ref}}$, and $T_{\mathrm{ref}}$ denote
fixed velocity, pressure, and time scales, respectively, and let
$|\Omega|=\boldsymbol{w}^{\mathsf T}\boldsymbol{1}$ denote the discrete
domain area.  These reference scales are common to all regime-local charts.
If the full-order equations and snapshots have already been
nondimensionalized, the corresponding reference scales are equal to one.

The dimensionless velocity and pressure energy measures at time $t_i$ are
defined by
\begin{equation}
  \widehat{E}^{u}_i
  =
  \frac{
    \|\boldsymbol{u}_i\|_{M_u}^{2}
  }{
    |\Omega|U_{\mathrm{ref}}^{2}
  },
  \qquad
  \widehat{E}^{p}_i
  =
  \frac{
    \|p_i^{\circ}\|_{M_p}^{2}
  }{
    |\Omega|P_{\mathrm{ref}}^{2}
  }.
  \label{eq:dimensionless_physical_energies}
\end{equation}
Here $\widehat{E}^{u}_i$ is the mean squared velocity magnitude and
$\widehat{E}^{p}_i$ is the corresponding mean squared magnitude of the
gauge-fixed pressure field.  These quantities measure the instantaneous
amplitudes of the two physical fields without reference to any local POD
basis.

The dimensionless first-order rates of change over the most recent time
interval are
\begin{equation}
  \widehat{D}^{u}_n
  =
  \frac{
    T_{\mathrm{ref}}
  }{
    \Delta t_n\sqrt{|\Omega|}U_{\mathrm{ref}}
  }
  \left\|
    \boldsymbol{u}_n-\boldsymbol{u}_{n-1}
  \right\|_{M_u},
  \qquad
  \widehat{D}^{p}_n
  =
  \frac{
    T_{\mathrm{ref}}
  }{
    \Delta t_n\sqrt{|\Omega|}P_{\mathrm{ref}}
  }
  \left\|
    p_n^{\circ}-p_{n-1}^{\circ}
  \right\|_{M_p}.
  \label{eq:dimensionless_physical_change_rates}
\end{equation}
They quantify the temporal variation of the velocity and pressure fields
relative to the prescribed physical scales.

To distinguish growth, decay, and saturation of the velocity amplitude, we
also introduce the dimensionless logarithmic energy-growth rate
\begin{equation}
  \widehat{g}_n
  =
  \frac{
    T_{\mathrm{ref}}
  }{
    \Delta t_n
  }
  \log
  \left(
    \frac{
      \widehat{E}^{u}_n+\varepsilon_E
    }{
      \widehat{E}^{u}_{n-1}+\varepsilon_E
    }
  \right),
  \label{eq:dimensionless_log_energy_growth}
\end{equation}
where $\varepsilon_E>0$ prevents numerical singularities when the measured
energy is close to zero.  In contrast to the absolute change rate
$\widehat{D}^{u}_n$, the quantity $\widehat{g}_n$ measures the relative
change of the velocity energy.

The three-state history is used to construct a dimensionless second-order
temporal-change indicator.  For nonuniform time increments, the discrete
velocity acceleration is approximated by
\begin{equation}
  \boldsymbol{A}^{u}_n
  =
  \frac{2}{
    \Delta t_n+\Delta t_{n-1}
  }
  \left[
    \frac{
      \boldsymbol{u}_n-\boldsymbol{u}_{n-1}
    }{
      \Delta t_n
    }
    -
    \frac{
      \boldsymbol{u}_{n-1}-\boldsymbol{u}_{n-2}
    }{
      \Delta t_{n-1}
    }
  \right].
  \label{eq:discrete_velocity_acceleration}
\end{equation}
The associated dimensionless temporal-curvature indicator is
\begin{equation}
  \widehat{\kappa}^{u}_n
  =
  \frac{
    T_{\mathrm{ref}}^{2}
  }{
    \sqrt{|\Omega|}U_{\mathrm{ref}}
  }
  \left\|
    \boldsymbol{A}^{u}_n
  \right\|_{M_u}.
  \label{eq:dimensionless_velocity_curvature}
\end{equation}
This quantity measures changes in the direction or magnitude of the recent
velocity increment and therefore complements the first-order rate
$\widehat{D}^{u}_n$.

The resulting dimensionless physical-statistics vector is
\begin{equation}
  \boldsymbol{s}_n
  =
  \left[
    \widehat{E}^{u}_n,\,
    \widehat{E}^{p}_n,\,
    \widehat{D}^{u}_n,\,
    \widehat{D}^{p}_n,\,
    \widehat{g}_n,\,
    \widehat{\kappa}^{u}_n
  \right]^{\mathsf T}
  \in\mathbb{R}^{6}.
  \label{eq:dimensionless_physical_statistics}
\end{equation}
Before being supplied to the overlap-fusion model, these quantities are
standardized componentwise using statistics computed exclusively from the
corresponding training population:
\begin{equation}
  \boldsymbol{d}_n
  =
  \left(
    \boldsymbol{s}_n-\boldsymbol{\mu}_{s}
  \right)
  \oslash
  \left(
    \boldsymbol{\sigma}_{s}
    +\varepsilon_s\boldsymbol{1}
  \right),
  \qquad
  \boldsymbol{d}_n\in\mathbb{R}^{6},
  \label{eq:standardized_physical_descriptors}
\end{equation}
where $\boldsymbol{\mu}_{s}$ and $\boldsymbol{\sigma}_{s}$ are the
componentwise training-set mean and standard deviation, respectively,
$\oslash$ denotes componentwise division, and $\varepsilon_s>0$ prevents
division by a vanishing standard deviation.  Validation and test
trajectories do not contribute to these normalization statistics.

The descriptor set in
\eqref{eq:dimensionless_physical_statistics} is intentionally restricted to
a small number of nonredundant quantities.  The two energy measures describe
the current field amplitudes, the two first-order rates characterize the
instantaneous temporal variation, the logarithmic growth rate distinguishes
relative growth from decay, and the second-order indicator captures changes
in the recent velocity evolution.  Previous-step rates and additional
nonlinear transformations are not included because they are strongly
correlated with these quantities and do not introduce a distinct category of
physical information.  The subsequent standardization places all six
components on comparable numerical scales and allows the fusion model to
determine their relative importance from the training data.

Since all components are computed from physical fields on the common mesh,
using the same quadrature matrices, pressure gauge, and reference scales,
$\boldsymbol{d}_n$ is independent of the regime-local POD basis used to
represent or advance the state.  It can therefore be evaluated consistently
for predictions initialized from either adjacent chart.



Let $\widehat{\boldsymbol{q}} = [\widetilde{\mu}; \boldsymbol{d}_n] \in \mathbb{R}^{7}$ 
denote the concatenation of the standardized Reynolds number 
$\widetilde{\mu}$ defined in \eqref{eq:e2_outer_router} and the 
standardized physical descriptors $\boldsymbol{d}_n \in \mathbb{R}^{6}$ 
defined in \eqref{eq:standardized_physical_descriptors}.
A lightweight correction network
$c_{\vartheta}$ (two hidden layers of width 64 with SiLU activations,
final linear layer zero-initialized) produces a scalar adjustment
$\delta = c_{\vartheta}(\widehat{\boldsymbol{q}})$.  The fusion weight is
then
\begin{equation}
  \alpha
  =
  \operatorname{sigmoid}
  \left(
  \ell^{\mathrm{E2}} + \delta
  \right),
  \qquad
  \alpha\in(0,1).
  \label{eq:t2c_convex_gate}
\end{equation}
Zero initialization guarantees that at the start of gate training
$\delta=0$ and the blend exactly reproduces the frozen E2 pair probability.
The complementary weight is $1-\alpha$, so the combination is strictly
convex.

The two specialists advance independently for $K$ steps from the same
three-state initial history.  Both use the unified frozen-context RK4
velocity update~\eqref{eq:unified_rk4} and the gated pressure
closure~\eqref{eq:adaptive_pressure_closure}; they differ only in their
regime-specific parameter sets $\theta_r$.  The velocity and
gauge-corrected pressure fields are decoded and then fused at each query
time $t_q$:
\begin{align}
  \widehat{\boldsymbol{u}}(t_q)
  &=
  \alpha\,\widehat{\boldsymbol{u}}_{1}(t_q)
  + (1-\alpha)\,\widehat{\boldsymbol{u}}_{2}(t_q),
  \label{eq:field_fusion}\\
  \widehat{p}^{\circ}(t_q)
  &=
  \alpha\,\widehat{p}^{\circ}_{1}(t_q)
  + (1-\alpha)\,\widehat{p}^{\circ}_{2}(t_q),
  \label{eq:pressure_field_fusion}
\end{align}
where the subscripts $1$ and $2$ refer to the candidate ordering of the
respective boundary.

The gate is trained on the training population and selected on validation
data using a sealed cache that stores the per-specialist prediction errors
and the corresponding quadratic statistics.  For a single field channel
$c\in\{u,p\}$, let $\boldsymbol{e}_1^c$ and $\boldsymbol{e}_2^c$ be the
errors of the two specialists with respect to the true field.  The squared
norm of the fused error is
\begin{equation}
  Q^c(\alpha)
  =
  \alpha^2\|\boldsymbol{e}_1^c\|_{M_c}^{2}
  +(1-\alpha)^2\|\boldsymbol{e}_2^c\|_{M_c}^{2}
  +2\alpha(1-\alpha)
  \langle\boldsymbol{e}_1^c,\boldsymbol{e}_2^c\rangle_{M_c}.
  \label{eq:t2c_quadratic_error}
\end{equation}
The training objective minimizes the average relative squared error over
all $K$ prediction steps and over the training windows:
\begin{equation}
  \mathcal{L}_{\mathrm{T2C}}
  =
  \operatorname{mean}_{w,k}
  \left[
  \frac{Q^u(\alpha^{(w)})}
       {\|\boldsymbol{u}_{\rm true}\|_{M_u}^{2}+\varepsilon}
  +
  \frac{Q^p(\alpha^{(w)})}
       {\|p_{\rm true}^{\circ}\|_{M_p}^{2}+\varepsilon}
  \right].
  \label{eq:t2c_training_loss}
\end{equation}
The gate checkpoint is selected by lexicographically minimizing
$(\max_{\mu}\ \mathrm{mean\ error},\ \mathrm{global\ mean\ error})$ on the
validation set, which prioritizes worst-case Reynolds-number robustness.
The fused output is never fed back into either specialist.

\paragraph{Application to the $S$--$H$ and $H$--$P$ boundaries}
For the $S$--$H$ boundary, the gate is trained on windows from the
steady-native development domain; Candidate~1 is the Steady specialist and
Candidate~2 is the Hopf specialist.  For the $H$--$P$ boundary, Candidate~1
is the Periodic specialist and Candidate~2 is the Hopf specialist.  In both
cases, the gate is trained and selected exclusively on the corresponding
training and validation windows, and the frozen specialists are evaluated
identically to their native rollouts.

Admissibility acts as a safety filter: if, for a given window, one of the
two specialists fails the admissibility check, the mask~\eqref{eq:masked_router_weights}
degenerates the prediction to the single admissible specialist.  In the
present validated configuration, both specialists passed the preflight
checks on all development and heldout windows, so the T2-C gate is always
active.  The $H$--$P$ boundary is thus treated as a full fusion boundary
rather than a hard fallback, consistent with the validated evidence that
both the periodic and Hopf specialists deliver stable native rollouts in
this region.

\paragraph{Complete physical-field predictor}
Let
$\Omega_S,\Omega_{SH},\Omega_H,\Omega_{HP},\Omega_P$
denote the five operating regions frozen from training and validation
evidence.  Combining the core, fusion, and fail-closed branches, the
physical-field predictor is
\begin{equation}
  \widehat{\boldsymbol{x}}^{\,\mathrm{phys}}
  =
  \begin{cases}
  \mathcal{D}_{S}\!\circ\mathcal{F}_{S},
  &(\mathcal{H},\mu)\in\Omega_S,\\[1mm]
  \alpha_{SH}\,\mathcal{D}_{S}\!\circ\mathcal{F}_{S}
  +(1-\alpha_{SH})\,\mathcal{D}_{H}\!\circ\mathcal{F}_{H},
  &(\mathcal{H},\mu)\in\Omega_{SH},\\[1mm]
  \mathcal{D}_{H}\!\circ\mathcal{F}_{H},
  &(\mathcal{H},\mu)\in\Omega_H,\\[1mm]
  \alpha_{HP}\,\mathcal{D}_{P}\!\circ\mathcal{F}_{P}
  +(1-\alpha_{HP})\,\mathcal{D}_{H}\!\circ\mathcal{F}_{H},
  &(\mathcal{H},\mu)\in\Omega_{HP},\\[1mm]
  \mathcal{D}_{P}\!\circ\mathcal{F}_{P},
  &(\mathcal{H},\mu)\in\Omega_P,\\[1mm]
  \mathrm{FAIL\mbox{-}CLOSED},
  &\sum_{j\in\mathcal{C}_{\mathrm{top}}}m_j=0.
  \end{cases}
  \label{eq:complete_ral_moe_predictor}
\end{equation}
In the fusion branches, $\alpha_{SH}$ and $\alpha_{HP}$ are the
window-conditioned convex weights produced by the respective T2-C gates.
When only one adjacent specialist is admissible, the corresponding fusion
branch reduces to the single admissible specialist by construction of the
masked weights.  The sum in the fail-closed branch is taken over the
topological candidate set; if no candidate is admissible, no numerical
prediction is issued.


















\subsection{Offline--online workflow and certified methodological scope}
\label{subsec:offline_online_scope}

\begin{algorithm}[H]
  \caption{Offline construction of RAL-MoE-ROM}
  \label{alg:ral_moe_offline}
  \KwIn{Complete Reynolds-number trajectories and registered native
  specialist contracts.}
  \KwOut{Frozen local specialists, outer router, admissibility records, and
  the certified $S$--$H$ and $H$--$P$ T2-C gates.}
  Partition complete trajectories into disjoint training, validation, and
  test Reynolds-number populations\;
  Construct train-only local velocity and pressure POD charts, means,
  scalers, Galerkin operators, and Pressure--Poisson assets\;
  Train the $S$, $H$, and $P$ HPRS-MoE-ROM specialists independently
  using the unified RK4 velocity update and gated pressure closure\;
  Select each specialist with validation trajectories only and freeze all
  local assets\;
  Train and freeze the E2 trajectory-level outer router\;
  Apply the adjacency topology in \eqref{eq:adjacency_matrix}\;
  Certify each candidate specialist on its registered development domain and
  horizon (admissibility preflight)\;
  Train T2-C gates on the validation-supported $S$--$H$ and $H$--$P$
  boundaries, using the respective chart-independent descriptors
  \eqref{eq:standardized_physical_descriptors}\;
  Freeze all thresholds, gates, and hierarchical decisions before final
  evaluation\;
\end{algorithm}

\begin{algorithm}[H]
  \caption{Hierarchical online prediction with RAL-MoE-ROM}
  \label{alg:ral_moe_online}
  \KwIn{$\mu$, a valid three-state physical history, query times, and a
  registered horizon $K$.}
  \KwOut{Physical velocity, gauge-corrected pressure, active specialists,
  weights, and admissibility diagnostics.}
  Evaluate the trajectory-level E2 router once\;
  restrict the candidates by the adjacency topology\;
  evaluate the admissibility mask for the candidate specialists\;
  \eIf{no candidate is admissible}{
    return fail-closed diagnostics without a forced prediction\;
  }{
    \eIf{exactly one candidate is admissible}{
      execute its frozen native rollout and decode the physical fields\;
    }{
      map the common physical history into the two admissible local charts\;
      execute two independent frozen native rollouts at aligned query times\;
      evaluate the corresponding T2-C gate once (trained for $S$--$H$ or
      $H$--$P$) and keep its weight $\alpha$ fixed over the window\;
      decode both predictions and fuse only the physical outputs by
      \eqref{eq:field_fusion}--\eqref{eq:pressure_field_fusion}\;
    }
  }
  return predictions and the complete routing/admissibility record\;
\end{algorithm}

Under Top‑1 routing, one frozen specialist is active.  The legal Top‑2
case evaluates two frozen native rollouts; the gate cost is small relative
to this second rollout.  No common POD decoder is needed and no specialist
parameter is updated online.  Exact wall‑clock acceleration factors are
not inferred without direct timing measurements.

The certified methodological scope can now be stated precisely.
RAL‑MoE‑ROM does not introduce a common public POD state, direct
cross‑chart modal blending, a full‑state continuous vector field, a
pressure derivative, cross‑specialist RHS fusion, fused‑state feedback, a
direct $S$--$P$ route, Top‑3 activation, or per‑step outer regime
switching.  All three specialists employ the unified frozen‑context RK4
velocity update~\eqref{eq:unified_rk4} and the gated pressure
closure~\eqref{eq:adaptive_pressure_closure}; no explicit Euler step,
additive pressure correction, or regime‑specific temporal encoding is
used.  The only soft combination occurs on the two adjacent boundaries
$S$--$H$ and $H$--$P$, where a window‑conditioned, output‑only convex
field fusion is applied via the T2‑C gates.  When admissibility flags a
specialist as unusable, the hierarchy falls back to the single admissible
candidate; if no candidate is admissible, the model fails closed.  These
restrictions define the verified model, rather than engineering omissions
to be silently replaced by stronger untested claims.



